\documentclass{article}
\usepackage[margin=1.15in]{geometry}
\usepackage{xcolor}
\usepackage{parskip}
\usepackage{fancyhdr}
\usepackage{array}
\usepackage{booktabs}
\usepackage{enumitem}
\usepackage{amsmath}
\usepackage{tabularx}
\usepackage{listings}

\pagestyle{fancy}
\fancyhf{}
\cfoot{\thepage}
\rhead{Lecture 14}
\lhead{MA 214: Applied Statistics (Spring 2026)}
\renewcommand{\headrulewidth}{.4mm}

\newcommand{\blankline}[1]{\underline{\hspace{#1}}}

\lstset{
  basicstyle=\ttfamily\small,
  columns=fullflexible,
  frame=single,
  framerule=0.4pt,
  breaklines=true,
  showstringspaces=false
}

\begin{document}

\noindent\textbf{Name:}\blankline{2.25in}\hfill\textbf{BUID:}\blankline{1.55in}

\medskip

\noindent\textbf{Name:}\blankline{2.25in}\hfill\textbf{BUID:}\blankline{1.55in}\newline

\begin{center}
 \Large In-Class Worksheet: Inference for Multiple Regression
\end{center}

\subsection*{Scenario}
The \emph{World Happiness Report} records $n = 150$ countries. We model \texttt{life\_expectancy} (years) from \texttt{gdp\_per\_capita} and \texttt{social\_support} (a $0$--$1$ index of whether people say they have someone to count on).

\smallskip
\textbf{Careful:} the \texttt{gdp\_per\_capita} column in this file is \emph{already on a log scale}. \texttt{exp(gdp\_per\_capita)} recovers GDP in dollars; we call that \texttt{gdp}.

\newpage

\subsection*{Part 1: Diagnose, then transform}

Two models for life expectancy, each with GDP as the only predictor:

\begin{center}
\renewcommand{\arraystretch}{1.3}
\begin{tabular}{lrrr}
\toprule
model & slope & $R^2$ & residual pattern \\
\midrule
\texttt{life\_expectancy $\sim$ gdp} (dollars)   & $0.00072$ & $0.429$ & clearly curved \\
\texttt{life\_expectancy $\sim$ log(gdp)}        & $3.81$    & $0.655$ & no curve \\
\bottomrule
\end{tabular}
\end{center}

\begin{enumerate}[label={\bfseries (\Alph*)},itemsep=.8em]
\item The raw-dollar model curves. Sketch, in words, what its residuals-vs-fitted plot must look like, and say which of the three misspecifications (non-linearity / non-constant variance / non-normality) this is. \\[3.5em]
\item Which variable did we transform, and why is that the right side here --- given that the Q--Q plot of the raw model was already roughly straight? \\[3.5em]
\item Interpret the slope $3.81$ of the $\log(\text{gdp})$ model: what happens to life expectancy when a country's GDP per capita \emph{doubles}? Show the arithmetic. \hfill \blankline{1.4in} years \\[2em]
\end{enumerate}

\subsection*{Part 2: Read the software output}

Now add social support, and run \texttt{tidy(fit, conf.int = TRUE)} on

\begin{center}
\texttt{lm(life\_expectancy $\sim$ gdp\_per\_capita + social\_support)}
\end{center}

\begin{center}
\renewcommand{\arraystretch}{1.3}
\begin{tabular}{lrrrrrr}
\toprule
term & estimate & std.error & statistic & p.value & conf.low & conf.high \\
\midrule
\texttt{(Intercept)}     & $38.50$ & $1.56$ & $24.71$ & $<0.001$ & $35.42$ & $41.58$ \\
\texttt{gdp\_per\_capita} & $2.98$  & $0.22$ & $13.48$ & $<0.001$ & $2.54$  & $3.41$ \\
\texttt{social\_support}  & $7.71$  & $1.00$ & $7.67$  & $<0.001$ & $5.73$  & $9.70$ \\
\bottomrule
\end{tabular}
\end{center}

\begin{enumerate}[label={\bfseries (\Alph*)},itemsep=.8em]
\item How many degrees of freedom does each $t$ statistic have? Give the number \emph{and} the reason. \hfill $\mathrm{df} = \blankline{1in}$ \\[1.5em]
\item Confirm the \texttt{social\_support} row: compute $t = \hat\beta / \mathrm{SE}$, and rebuild its interval from $\hat\beta \pm t^\ast\,\mathrm{SE}$ using $t^\ast = 1.976$. \\[3.5em]
\item Interpret $\hat\beta_{\texttt{social\_support}} = 7.71$ in a full sentence --- with units, and with the phrase this lecture insisted on. \\[3.5em]
\item The GDP slope \emph{fell} from $3.81$ (Part 1, GDP alone) to $2.98$ once social support joined the model. Explain why that is not a contradiction --- what changed about the question being asked? \\[4em]
\end{enumerate}

\newpage

\subsection*{Part 3: Find the bugs}

A classmate wants the effect of social support on life expectancy, adjusted for GDP, with a $95\%$ interval. Their code runs without a single error message --- and is wrong in \textbf{three} separate ways.

\begin{lstlisting}[language=R]
happiness <- read_csv("world_happiness.csv") |>
  mutate(gdp = exp(gdp_per_capita))   # GDP in dollars

# The effect of social support, holding GDP fixed:
fit1 <- lm(life_expectancy ~ social_support, data = happiness)
tidy(fit1)

# Same thing, now with GDP on the log scale:
fit2 <- lm(life_expectancy ~ log(gdp_per_capita) + social_support,
           data = happiness)

# 95% CI for the social_support slope
b  <- coef(fit2)["social_support"]
se <- coef(summary(fit2))["social_support", "Std. Error"]
b + c(-1, 1) * qt(0.975, df = nrow(happiness) - 2) * se
\end{lstlisting}

\begin{enumerate}[label={\bfseries (\Alph*)},itemsep=.8em]
\item \textbf{Bug 1} is the comment above \texttt{fit1} --- the code is fine, the \emph{claim} is not. What does the \texttt{fit1} coefficient actually measure, and why can it not be read ``holding GDP fixed''? \\[4em]
\item \textbf{Bug 2} is inside \texttt{fit2}'s formula. Look again at what the \texttt{gdp\_per\_capita} column already is. What has the code done to it, and what should the term be? \\[4em]
\item \textbf{Bug 3} is the \texttt{df} in the last line. What should it be, and why? \hfill $\mathrm{df} = \blankline{1in}$ \\[2.5em]
\end{enumerate}

\subsection*{Part 4: Does the effect of GDP depend on social support?}

Fit the interaction, \texttt{lm(life\_expectancy $\sim$ gdp\_per\_capita * social\_support)}:

\begin{center}
\renewcommand{\arraystretch}{1.3}
\begin{tabular}{lrrrr}
\toprule
term & estimate & std.error & statistic & p.value \\
\midrule
\texttt{(Intercept)}                      & $37.00$  & $3.11$ & $11.90$ & $<0.001$ \\
\texttt{gdp\_per\_capita}                  & $3.18$   & $0.42$ & $7.55$  & $<0.001$ \\
\texttt{social\_support}                   & $11.58$  & $6.89$ & $1.68$  & $0.096$ \\
\texttt{gdp\_per\_capita:social\_support}   & $-0.491$ & $0.87$ & $-0.56$ & $0.574$ \\
\bottomrule
\end{tabular}
\end{center}

\smallskip
\noindent Model comparison, all on the same $150$ countries:

\begin{center}
\begin{tabular}{lrr}
\toprule
model & adj.\ $R^2$ & AIC \\
\midrule
\texttt{gdp\_per\_capita + social\_support} & $0.750$ & $708.5$ \\
\texttt{gdp\_per\_capita * social\_support} & $0.749$ & $710.2$ \\
\bottomrule
\end{tabular}
\end{center}

\begin{enumerate}[label={\bfseries (\Alph*)},itemsep=.8em]
\item Write down the slope of \texttt{gdp\_per\_capita} as a function of \texttt{social\_support}, using the fitted numbers. \\[2.5em]
\item Evaluate that slope for a country with weak social support ($0.7$) and one with strong social support ($0.95$). Are the two answers far apart? \\[3.5em]
\item Should we keep the interaction? Cite the $p$-value, the AIC, \emph{and} (C) --- and say what the answer is in plain English, about GDP and social support in the real world. \\[4.5em]
\item \texttt{social\_support} was overwhelming in Part 2 ($p < 0.001$) and is now unconvincing ($p = 0.096$). Nothing about the data changed. What happened --- and which idea from today explains it? \\[4.5em]
\end{enumerate}

\newpage

\subsection*{Part 5: Build your own model}

You have \texttt{freedom}, \texttt{generosity}, \texttt{corruption\_perception} and \texttt{unemployment\_rate} still unused.

\begin{enumerate}[label={\bfseries (\Alph*)},itemsep=.8em]
\item Add \emph{one} of them to the Part 2 model. Report its coefficient, its $p$-value, and what happens to the other two coefficients. \\[4em]
\item Check the diagnostics of your model (residuals vs fitted, Q--Q, and residuals against each predictor). Does anything need transforming? \\[4em]
\item Compute the correlations among your predictors. Is any pair collinear enough to worry about? What would that do to your table --- and would it damage a \emph{prediction} of life expectancy for a new country? \\[4.5em]
\item Which model would you actually report, and why? One paragraph. Do not say ``the one with the highest $R^2$'' unless you can defend it. \\[5em]
\end{enumerate}

\subsection*{Compare with the groups around you}

\noindent\fbox{\parbox{\dimexpr\textwidth-2\fboxsep-2\fboxrule\relax}{%
Compare your final model with at least two neighboring pairs. Did you keep the same predictors? The same transformation? If not --- was it a different \emph{judgment}, or a different \emph{mistake}? Having heard them, would you change your model?
\\[6em]
}}

\end{document}
